%% toc: P1 1 1.5
\section{Changing the formal language of a logical system}
\label{sec:p01-w02-p01-003-language-change}

A logical system \(\langle\Lang,\mathcal C,\mathcal S\rangle\) has three
components, and the example after the definition in
\cref{sec:levels-01} varies only two of them: it replaces a
natural-deduction calculus by a sequent calculus, and two-valued valuations
by three-valued matrices, while the propositional language stays fixed.
This section treats the remaining case, a change of the language itself.
Its purpose is to say exactly what has to be specified again when
\(\Lang\) changes, and exactly where the old and the new system can then be
compared.

The case is the first-order one. The input is a first-order signature
\(\Lang\), a second signature \(\Lang'\) extending it, the calculus of
first-order natural deduction and the semantics of first-order structures
with variable assignments, all as defined in the inherited first-order
chapters \cite[chs.~``First-Order Syntax'', ``Structures, Satisfaction, and
Translation'', ``First-Order Deduction'']{FrozenFormalLogicLanguageChange}.
The conclusions are a re-specification procedure
(\cref{def:p01-w02-p01-003-respecification}), two comparison theorems for the
formulas of the smaller language
(\cref{thm:p01-w02-p01-003-semantic-comparison,thm:p01-w02-p01-003-derivability-comparison}),
a comparison through a common sub-signature
(\cref{thm:p01-w02-p01-003-common-subsignature}), and a case in which the
comparison fails because the class of permitted structures was narrowed at the
same time (\cref{thm:p01-w02-p01-003-restricted-class}). The arguments are
editorial completions assembled from the inherited definitions and results;
they are not proofs attributed to an external presentation.

\subsection{The logical system over a signature}
\label{sec:p01-w02-p01-003-system-over-signature}

A first-order signature
\(\Lang=\langle\mathcal F,\mathcal R,\operatorname{ar}_{\mathcal F},
\operatorname{ar}_{\mathcal R}\rangle\) lists function symbols, relation
symbols and their arities. A signature \(\Lang'\) \emph{extends}
\(\Lang\) when \(\mathcal F\subseteq\mathcal F'\),
\(\mathcal R\subseteq\mathcal R'\) and the arity maps of \(\Lang'\) restrict
to those of \(\Lang\). For an \(\Lang'\)-structure \(\mathcal M'\), the
\emph{reduct} \(\mathcal M'\mathbin{\upharpoonright}\Lang\) keeps the domain
and the interpretations of the symbols of \(\Lang\) and forgets the others;
an \(\Lang'\)-structure is an \emph{expansion} of an \(\Lang\)-structure
\(\mathcal M\) when it has the domain of \(\mathcal M\) and agrees with
\(\mathcal M\) on every symbol of \(\Lang\)
\cite[ch.~``Structures, Satisfaction, and Translation'', definition
``Signature extension, expansion, and reduct'']{FrozenFormalLogicLanguageChange}.
Every structure has a nonempty domain \cite[ch.~``Structures, Satisfaction,
and Translation'', definition ``First-order structure'']{FrozenFormalLogicLanguageChange}.

The inherited definitions of the formula set, of semantic consequence and of
first-order derivability are each stated for one signature fixed in the
surrounding text: \(\Gamma\entails\varphi\) quantifies over the
\(\Lang\)-structures, and the rule instances of the calculus range over
\(\Lang\)-formulas and \(\Lang\)-terms. Once two signatures are active, the
signature has to be written into the notation, as \cref{sec:levels-01}
requires whenever more than one semantics or calculus is in play.

\begin{hbconvention}{def:p01-w02-p01-003-signature-subscripts}{Consequence and derivability over a named signature}
For a signature \(\Lang\), write
\(\operatorname{Term}(\Lang)\), \(\operatorname{Fm}(\Lang)\) and
\(\operatorname{Sent}(\Lang)\) for its terms, formulas and sentences.
For \(\Gamma\cup\{\varphi\}\subseteq\operatorname{Fm}(\Lang)\):
\begin{enumerate}
  \item \(\Gamma\entails_{\Lang}\varphi\) means that for every
        \(\Lang\)-structure \(\mathcal M\) and every assignment \(s\) in its
        domain, if \(\mathcal M,s\entails\gamma\) for every
        \(\gamma\in\Gamma\), then \(\mathcal M,s\entails\varphi\);
  \item \(\mathsf{ND}_{\mathsf{FOL}}[\Lang]\) is the calculus whose rules are
        the rule schemata of first-order natural deduction (those of
        \(\mathsf{ND}_{\mathsf{CPL}}\), the quantifier rules
        \(\forall E\), \(\exists I\), \(\forall I\), \(\exists E\), the
        conversion rule \(\alpha\), and the identity rules \(=I\), \(=E\)),
        with every formula in a rule instance drawn from
        \(\operatorname{Fm}(\Lang)\) and every instantiating term drawn from
        \(\operatorname{Term}(\Lang)\); and
        \(\Gamma\derives_{\mathsf{ND}_{\mathsf{FOL}}[\Lang]}\varphi\) means
        that some derivation in this calculus has conclusion \(\varphi\) and
        all its open assumptions in \(\Gamma\).
\end{enumerate}
The logical system over \(\Lang\) is
\(\langle\operatorname{Fm}(\Lang),\mathsf{ND}_{\mathsf{FOL}}[\Lang],
\mathcal S_{\Lang}\rangle\), where \(\mathcal S_{\Lang}\) permits as semantic
data exactly the pairs \((\mathcal M,s)\) of an \(\Lang\)-structure and an
assignment in it. With one signature fixed, the subscripts may be dropped
and the inherited notation \(\Gamma\entails\varphi\) and
\(\Gamma\derives_{\mathsf{ND}_{\mathsf{FOL}}}\varphi\) results
\cite[ch.~``Structures, Satisfaction, and Translation'', definition
``Truth, validity, and semantic consequence''; ch.~``First-Order
Deduction'', definition ``First-order derivability'']{FrozenFormalLogicLanguageChange}.
\index{logical system!change of language}\index{consequence!over a named signature}
\end{hbconvention}

\subsection{What must be specified again}
\label{sec:p01-w02-p01-003-respecify}

\begin{hbmethod}{def:p01-w02-p01-003-respecification}{Re-specifying a first-order system after a change of signature}
Let \(\Lang'\) extend \(\Lang\). Passing from the system over \(\Lang\) to
the system over \(\Lang'\) requires the following five specifications, none
of which is inherited automatically from the old system.
\begin{enumerate}
  \item \emph{Grammar.} \(\operatorname{Term}(\Lang')\) and
        \(\operatorname{Fm}(\Lang')\) are generated by the same clauses with
        the symbols of \(\Lang'\). Side condition: the new symbols must be
        disjoint from the variables, from the logical and punctuation tokens
        and from each other, and each must receive one arity; the signature
        definition requires exactly this \cite[ch.~``First-Order Syntax'',
        definition ``First-order signature'']{FrozenFormalLogicLanguageChange}.
        The variable set \(\mathsf{Var}\) and its fixed order are unchanged.
  \item \emph{Semantic data.} The permitted data become pairs
        \((\mathcal M',s)\) of an \(\Lang'\)-structure and an assignment.
        Every new function symbol needs a total function on the domain of
        the matching arity, every new constant an element of the domain,
        every new relation symbol a relation of the matching arity. The
        satisfaction clauses are applied to the larger formula set; for a new
        relation symbol the atomic clause reads off its new interpretation.
  \item \emph{Calculus.} The rule schemata stay as they are, but their
        instances now range over \(\operatorname{Fm}(\Lang')\) and
        \(\operatorname{Term}(\Lang')\). In particular \(\forall E\),
        \(\exists I\), \(=I\) and \(=E\) may now use terms built from new
        symbols. The side conditions of \(\forall I\) and \(\exists E\) and
        the side condition of the rule \(\alpha\) concern variables only, so
        their wording does not change, but they are now checked on
        \(\Lang'\)-formulas.
  \item \emph{Consequence relations.} Both
        \(\entails_{\Lang'}\) and
        \(\derives_{\mathsf{ND}_{\mathsf{FOL}}[\Lang']}\) are new relations,
        defined on sets of \(\Lang'\)-formulas. That the old calculus was
        sound and complete for the old semantics says nothing about the new
        pair; soundness and completeness of
        \(\mathsf{ND}_{\mathsf{FOL}}[\Lang']\) for \(\mathcal S_{\Lang'}\) are
        the inherited theorems applied to the signature \(\Lang'\)
        (\cref{thm:fol-soundness,thm:fol-completeness}), and the second
        carries its Axiom-of-Choice hypothesis.
  \item \emph{Comparison domain.} The two systems can be compared only on
        premise sets and conclusions that are formulas of both languages.
        For an extension this is \(\operatorname{Fm}(\Lang)\), which lies
        inside \(\operatorname{Fm}(\Lang')\) by
        \cref{thm:p01-w02-p01-003-syntax-inclusion}. A claim
        \(\Gamma\entails_{\Lang}\varphi\) with
        \(\varphi\in\operatorname{Fm}(\Lang')\setminus\operatorname{Fm}(\Lang)\)
        is not false; it is not defined.
\end{enumerate}
\index{language change!re-specification}
\end{hbmethod}

\subsection{Comparison on the formulas of the smaller language}
\label{sec:p01-w02-p01-003-comparison}

\begin{hblemma}{thm:p01-w02-p01-003-syntax-inclusion}{proved=annex:p01-w02-p01-003-syntax-inclusion}{The old syntax inside the new}
Let \(\Lang'\) extend \(\Lang\).
\begin{enumerate}
  \item \(\operatorname{Term}(\Lang)\subseteq\operatorname{Term}(\Lang')\)
        and \(\operatorname{Fm}(\Lang)\subseteq\operatorname{Fm}(\Lang')\).
  \item For \(\varphi\in\operatorname{Fm}(\Lang)\),
        \(t\in\operatorname{Term}(\Lang)\) and \(x\in\mathsf{Var}\), the
        capture-avoiding substitution \(\varphi[t/x]\) computed over
        \(\Lang\) is the same string as the one computed over \(\Lang'\), and
        it lies in \(\operatorname{Fm}(\Lang)\). The occurrence sets
        \(\operatorname{Var}(\varphi)\) and free-variable sets
        \(\operatorname{FV}(\varphi)\) are likewise the same over both
        signatures.
  \item If \(\varphi,\psi\in\operatorname{Fm}(\Lang)\) are
        alpha-equivalent as \(\Lang\)-formulas, they are alpha-equivalent as
        \(\Lang'\)-formulas.
\end{enumerate}
\end{hblemma}

\begin{hblemma}{thm:p01-w02-p01-003-expansion}{proved=annex:p01-w02-p01-003-expansion}{Every structure of the old signature has an expansion}
Let \(\Lang'\) extend \(\Lang\). Every \(\Lang\)-structure \(\mathcal M\)
has an \(\Lang'\)-expansion \(\mathcal M'\), and every such expansion
satisfies \(\mathcal M'\mathbin{\upharpoonright}\Lang=\mathcal M\). The
expansion is obtained from one element of the domain of \(\mathcal M\) and
uses no choice principle.
\end{hblemma}

\begin{hbtheorem}{thm:p01-w02-p01-003-semantic-comparison}{proved=annex:p01-w02-p01-003-semantic}{Semantic consequence is unchanged on the old formulas}
Let \(\Lang'\) extend \(\Lang\). For every
\(\Gamma\cup\{\varphi\}\subseteq\operatorname{Fm}(\Lang)\),
\[
 \Gamma\entails_{\Lang}\varphi
 \quad\Longleftrightarrow\quad
 \Gamma\entails_{\Lang'}\varphi .
\]
In particular an \(\Lang\)-formula is valid over \(\Lang\)-structures iff
it is valid over \(\Lang'\)-structures.
\end{hbtheorem}

The left-to-right direction uses only that every \(\Lang'\)-structure has a
reduct to \(\Lang\) on which the old formulas keep their truth values
(\cref{lem:reduct-invariance}). The right-to-left direction also uses that
every \(\Lang\)-structure arises as such a reduct
(\cref{thm:p01-w02-p01-003-expansion}). The second requirement is the one that
fails in \cref{thm:p01-w02-p01-003-restricted-class}.

\begin{hbtheorem}{thm:p01-w02-p01-003-derivability-comparison}{proved=annex:p01-w02-p01-003-derivability}{Derivability on the old formulas}
Let \(\Lang'\) extend \(\Lang\), and let
\(\Gamma\cup\{\varphi\}\subseteq\operatorname{Fm}(\Lang)\).
\begin{enumerate}
  \item If \(\Gamma\derives_{\mathsf{ND}_{\mathsf{FOL}}[\Lang]}\varphi\),
        then \(\Gamma\derives_{\mathsf{ND}_{\mathsf{FOL}}[\Lang']}\varphi\);
        indeed every derivation in \(\mathsf{ND}_{\mathsf{FOL}}[\Lang]\) is,
        unchanged, a derivation in \(\mathsf{ND}_{\mathsf{FOL}}[\Lang']\)
        with the same conclusion and the same open assumptions.
  \item Assume the Axiom of Choice in the ambient set theory. If
        \(\Gamma\derives_{\mathsf{ND}_{\mathsf{FOL}}[\Lang']}\varphi\), then
        \(\Gamma\derives_{\mathsf{ND}_{\mathsf{FOL}}[\Lang]}\varphi\).
\end{enumerate}
\end{hbtheorem}

Clause~1 is a fact about derivations and holds without any semantic
input. Clause~2 is obtained here by passing through the semantics: soundness
over \(\Lang'\), then \cref{thm:p01-w02-p01-003-semantic-comparison}, then
completeness over \(\Lang\). It therefore carries the Axiom-of-Choice
hypothesis of the inherited completeness theorem
(\cref{thm:fol-completeness}). A derivation over \(\Lang'\) of an
\(\Lang\)-formula may pass through formulas containing new symbols; this
section does not give a syntactic transformation that removes them, and so
does not establish clause~2 without that hypothesis.

\begin{hbcorollary}{thm:p01-w02-p01-003-common-subsignature}{proved=annex:p01-w02-p01-003-shared}{Comparison through a common sub-signature}
Let \(\Lang_1\) and \(\Lang_2\) be signatures, neither required to extend
the other. Let \(\Lang_0\) have as function symbols those
\(f\in\mathcal F_1\cap\mathcal F_2\) with
\(\operatorname{ar}_{\mathcal F_1}(f)=\operatorname{ar}_{\mathcal F_2}(f)\),
as relation symbols those \(R\in\mathcal R_1\cap\mathcal R_2\) with
\(\operatorname{ar}_{\mathcal R_1}(R)=\operatorname{ar}_{\mathcal R_2}(R)\),
and the common arities. Then \(\Lang_0\) is a signature, both \(\Lang_1\)
and \(\Lang_2\) extend it, and for every
\(\Gamma\cup\{\varphi\}\subseteq\operatorname{Fm}(\Lang_0)\),
\[
 \Gamma\entails_{\Lang_1}\varphi
 \quad\Longleftrightarrow\quad
 \Gamma\entails_{\Lang_0}\varphi
 \quad\Longleftrightarrow\quad
 \Gamma\entails_{\Lang_2}\varphi .
\]
Under the Axiom of Choice the same chain holds for
\(\derives_{\mathsf{ND}_{\mathsf{FOL}}[\Lang_1]}\),
\(\derives_{\mathsf{ND}_{\mathsf{FOL}}[\Lang_0]}\) and
\(\derives_{\mathsf{ND}_{\mathsf{FOL}}[\Lang_2]}\).
\end{hbcorollary}

The corollary fixes \(\operatorname{Fm}(\Lang_0)\) as a domain on which two
unrelated signature changes can be compared. It does not determine whether
some string belongs to \(\operatorname{Fm}(\Lang_1)\cap\operatorname{Fm}(\Lang_2)\)
without belonging to \(\operatorname{Fm}(\Lang_0)\); that question needs the
unique-readability analysis of the first-order grammar and is not settled
here.

\subsection{Worked example: adding a constant}
\label{sec:p01-w02-p01-003-example}

\begin{hbexample}{def:p01-w02-p01-003-add-constant}{Adding a constant to a signature with one binary relation}
Let \(\Lang\) have no function symbols and one binary relation symbol
\(R\), and let \(\Lang'\) add one constant symbol \(c\) (a function symbol
of arity zero). Then \(\Lang'\) extends \(\Lang\).

\emph{Grammar.} \(\operatorname{Term}(\Lang)=\mathsf{Var}\) and
\(\operatorname{Term}(\Lang')=\mathsf{Var}\cup\{c\}\). The atomic
\(\Lang'\)-formulas include \(x=c\), \(c=c\), \(R(c,x)\), \(R(x,c)\) and
\(R(c,c)\), none of which is an \(\Lang\)-formula.

\emph{Semantic data.} An \(\Lang\)-structure is a nonempty set with a binary
relation on it; an \(\Lang'\)-structure adds a designated element
\(c^{\mathcal M'}\).

\emph{Calculus.} From the assumption \(\forall x\,\neg R(x,x)\), the rule
\(\forall E\) with the term \(c\) gives \(\neg R(c,c)\) in
\(\mathsf{ND}_{\mathsf{FOL}}[\Lang']\), so
\(\{\forall x\,\neg R(x,x)\}\derives_{\mathsf{ND}_{\mathsf{FOL}}[\Lang']}
\neg R(c,c)\). This one-step tree is not a derivation in
\(\mathsf{ND}_{\mathsf{FOL}}[\Lang]\): its conclusion is not an
\(\Lang\)-formula and its instantiating term is not an \(\Lang\)-term.

\emph{Outside the comparison domain.} The new system relates
\(\forall x\,\neg R(x,x)\) to \(\neg R(c,c)\); the old system has no
corresponding claim. The \(\Lang\)-sentence \(\exists x\,\neg R(x,x)\)
follows from \(\forall x\,\neg R(x,x)\) in both systems, because every
domain is nonempty, but it is a different formula, and nothing in either
system identifies it with \(\neg R(c,c)\).

\emph{Inside the comparison domain.} Put
\(\Gamma=\{\forall x\,\forall y\,(R(x,y)\to R(y,x))\}\) and
\(\varphi=\forall x\,R(x,x)\), both in \(\operatorname{Fm}(\Lang)\).
Let \(\mathcal N\) have domain \(\{0,1\}\) and
\(R^{\mathcal N}=\varnothing\). The premise holds in \(\mathcal N\) because
no pair is in \(R^{\mathcal N}\), and \(\varphi\) fails at \(0\). So
\(\Gamma\not\entails_{\Lang}\varphi\).
\Cref{thm:p01-w02-p01-003-semantic-comparison} then gives
\(\Gamma\not\entails_{\Lang'}\varphi\); a witness is the expansion of
\(\mathcal N\) with \(c^{\mathcal N'}=0\), which agrees with \(\mathcal N\)
on every \(\Lang\)-formula by \cref{lem:reduct-invariance}.
\index{language change!worked example}
\end{hbexample}

\subsection{Failure when the permitted class is narrowed}
\label{sec:p01-w02-p01-003-failure}

The comparison theorems assume that the new semantics permits every
\(\Lang'\)-structure. A language change is often made together with new
axioms about the new symbols, and the semantics is then the class of
\(\Lang'\)-structures that satisfy them. That is also a change of the
semantic class, and the right-to-left direction of
\cref{thm:p01-w02-p01-003-semantic-comparison} can then fail.

\begin{hblemma}{thm:p01-w02-p01-003-restricted-class}{proved=annex:p01-w02-p01-003-restricted}{A new constant with an axiom that is not conservative}
Let \(\Lang\) and \(\Lang'\) be as in
\cref{def:p01-w02-p01-003-add-constant}, and put
\begin{gather*}
 A\defeq\forall x\,(x=c\lor R(c,x))\in\operatorname{Sent}(\Lang'),\\
 \theta\defeq\exists y\,\forall x\,(x=y\lor R(y,x))\in\operatorname{Sent}(\Lang),
\end{gather*}
where \(x\) and \(y\) are distinct variables. Let \(\mathcal K_A\) be the
class of \(\Lang'\)-structures \(\mathcal M'\) with
\(\mathcal M'\entails A\), and for
\(\Delta\cup\{\chi\}\subseteq\operatorname{Fm}(\Lang')\) write
\(\Delta\entails_{\mathcal K_A}\chi\) when every member of
\(\mathcal K_A\) and every assignment in it that satisfy all of \(\Delta\)
satisfy \(\chi\). Then:
\begin{enumerate}
  \item \(\varnothing\entails_{\mathcal K_A}\theta\);
  \item \(\varnothing\not\entails_{\Lang}\theta\), witnessed by the
        structure \(\mathcal N\) of \cref{def:p01-w02-p01-003-add-constant};
  \item \(\mathcal N\) has no \(\Lang'\)-expansion in \(\mathcal K_A\); and
  \item \(\{A\}\derives_{\mathsf{ND}_{\mathsf{FOL}}[\Lang']}\theta\), by one
        application of \(\exists I\) with the term \(c\).
\end{enumerate}
\end{hblemma}

Clause~3 locates the failure: the step of
\cref{thm:p01-w02-p01-003-semantic-comparison} that lifts an
\(\Lang\)-structure to an permitted \(\Lang'\)-structure is no longer
available. Clauses~1 and~4 show the same effect on the semantic and on the
calculus side: the \(\Lang\)-consequences of the new system include
\(\theta\), which is not an \(\Lang\)-validity. When the added axiom is
instead an explicit definition of a new relation symbol by an
\(\Lang\)-formula, every model of the old theory has exactly one permitted
expansion, and the old consequences are preserved exactly
(\cref{thm:definitional-conservativity}).

\begin{hbwarning}{def:p01-w02-p01-003-boundary}{Comparison boundary for a change of language}
\begin{enumerate}
  \item The comparison of
        \cref{thm:p01-w02-p01-003-semantic-comparison,thm:p01-w02-p01-003-derivability-comparison}
        is made only on premise sets and conclusions in
        \(\operatorname{Fm}(\Lang)\). A formula with a new symbol has no
        status in the old system. Comparing it with an old formula needs a
        separately specified translation, which is a further piece of data
        and not part of either system.
  \item The comparison assumes that the semantics of the new system permits
        every \(\Lang'\)-structure together with every assignment. If the
        permitted class is narrowed, as in
        \cref{thm:p01-w02-p01-003-restricted-class}, conservativity must be
        proved for that class; it does not follow from the language change.
  \item The theorems concern signature extensions of classical first-order
        logic, in which the satisfaction clauses for the logical symbols are
        the same in both systems. A change of the logical vocabulary itself,
        for example replacing the primitive connectives or adding a modal
        operator, changes those clauses; neither theorem applies to it, and
        its comparison boundary has to be established for the systems in
        question.
  \item Clause~2 of \cref{thm:p01-w02-p01-003-derivability-comparison} is
        proved here only under the Axiom of Choice, through completeness.
\end{enumerate}
\index{language change!comparison boundary}
\end{hbwarning}

\subsection{Relationships and sources}
\label{sec:p01-w02-p01-003-relationships}

This section completes the example after the definition of a logical system
in \cref{sec:levels-01}, which changes \(\mathcal C\) and \(\mathcal S\)
with \(\Lang\) fixed. The definitions of signature, term, formula,
capture-avoiding substitution and alpha-equivalence are those of
\cite[ch.~``First-Order Syntax'']{FrozenFormalLogicLanguageChange}; structures,
expansions, reducts, satisfaction and semantic consequence are those of
\cite[ch.~``Structures, Satisfaction, and
Translation'']{FrozenFormalLogicLanguageChange}; the natural-deduction rules and
their side conditions are those of \cite[ch.~``First-Order Deduction'',
section ``Natural-deduction rules for quantifiers and
equality'']{FrozenFormalLogicLanguageChange}. The proofs use invariance under
reduct (\cref{lem:reduct-invariance}), soundness
(\cref{thm:fol-soundness}) and strong completeness
(\cref{thm:fol-completeness}). Conservativity of definitional extensions
(\cref{thm:definitional-conservativity}) is the theory-level case in which a
language change comes with a new axiom and the old consequences survive;
\cref{thm:p01-w02-p01-003-restricted-class} is a case in which they do not.
The compactness proof (\cref{thm:fol-compactness}) uses a language change of
the same kind, passing to a signature with added witness constants and back
by reduct.
